\begin{afterword}
\summaryhead

The post backs the previous day's claim that science knows exactly what evolution can do, by calculating how slow natural selection is. A gene with a 3 per cent advantage takes about 768 generations to spread through a population of 100,000, by the formula $2\ln(N)/s$, and has only about a 6 per cent chance of spreading at all, since the chance is about twice the advantage. Genes must spread by their bearers' own reproduction, with no investors or imitators to speed them up.

Complex adaptations take far longer, because each dependent part gains little until the part it depends on is common. And one species' invention does not help another. A human programmer, by contrast, can design a hundred interdependent parts in an afternoon, plan ahead, change parts together and share the result at once. Technology already beats biology in many ways. So natural selection is a poor model for human design, and should not be praised beyond what it deserves.

\responsehead

The post's arithmetic is right, and its main point holds. The figures agree, to within about 5 per cent, with the standard results, Kimura and Ohta's time to fixation and Haldane's 2s, and the worked numbers (768, 875, 2763 generations, 6 per cent, 0.1 per cent) recompute correctly. The contrast with a human designer who plans ahead and changes several parts at once is fair.

The post says evolution is simple enough to calculate ``exactly how stupid'' it is. The formulas are approximations for an idealized population of constant size, with one gene, a fixed advantage and Poisson-distributed offspring; with an effective population smaller than the census, even the 2s result changes. That does not hurt the argument, which needs only orders of magnitude. It does qualify the promise of exactness, and the previous day's claim that science has ``a very exact idea'' of evolution's capabilities.

The timetable for complex adaptations is the slow case. It assumes that each dependent part appears as a new mutation after the previous part has spread (``Only then can another allele B \ldots\ begin rising''). A part already present at low frequency starts to rise as soon as what it depends on is common, so the steps can overlap, and adaptation from such standing variation is expected to be faster.

The claim that other evolutions do not share an invention is true of snakes and foxes, the post's examples, but not of evolution in general. Bacteria pass genes between species, and antibiotic resistance spreads this way. The programmer who posts code to the Internet has a counterpart in biology, if not in the animals the post discusses.

In short: correct arithmetic for the simplest model of selection, offered as exact, with a timetable that assumes the slower case and a claim that one species' inventions do not pass to another, which does not hold for bacteria.
\end{afterword}
